The descriptive differences between single-electrode and full-montage operation were next tested for every electrode using session-paired comparisons. For each of the electrodes, that electrode's session-level $F_1$ operating alone was paired against the same electrode's session-level $F_1$ scored inside the full-montage run, on the sessions common to both. For example, Fp2 operating as the entire detector was compared with Fp2 scored inside the full-montage Proposed-Med detector, and likewise for every other electrode.
Each electrode was tested with a two-tailed Wilcoxon signed-rank test, paired by session, with $p$ Bonferroni-corrected over the electrodes tested within each dataset
%(separate families for Internal and Cao2018, the same electrode set as Figure~\ref{fig:region_performance}) and the matched-pairs rank-biserial correlation reported as the effect size,
%using the sign convention $\Delta F_1 = F_{1,\mathrm{single}} - F_{1,\mathrm{full}}$ throughout, so a negative value indicates lower $F_1$ under single-electrode operation.
The Figure~\ref{fig:exp1_single_channel_stats} shows, for every electrode, its session-mean $F_1$ operating alone overlaid on the same electrode's session-mean $F_1$ under the full montage, together with the adjusted-significance marker.

On Internal (Figure~\ref{fig:exp1_single_channel_stats}A), 14 of the 24 electrodes tested remained significantly different after correction, and every significant difference was a reduction under single-electrode operation, with $\Delta F_1$ ranging from $-7.02$ to $-0.15$ percentage among all 24 electrodes.
Despite retaining the two highest absolute single-electrode scores, neither Fp1 nor Fp2 reached significance after correction (adjusted $p=1.000$ for both), wheras the AF4 and AF3 did reach significance which is lower in the single-electode relative to the full montage configuraion.
Beyond the frontal electrodes, the Internal C3--C4 asymmetry noted above coexisted with a significant configuration effect at both electrodes, and most occipital electrodes reached significance different except O1 and P4, despite all having very low absolute $F_1$.

On Cao2018 (Figure~\ref{fig:exp1_single_channel_stats}B), 16 of the 18 electrodes tested remained significantly different after correction, again all reductions. Similarly, only the frontopolar pair, Fp1 and Fp2, did not reach significance after correction (adjusted $p=1.000$ for both), and F3 and F4, the third- and fourth-highest single-electrode scores, showed the two largest absolute reductions ($\Delta F_1 = -10.21$ and $-8.96$ percentage points, respectively).
Every central, parietal and occipital electrode reached significance, including both C3 and C4 (consistent with the C3/C4 symmetry already noted).
